\section{Carta de los Presidentes}

\begin{chairletter}{Jimena Solórzano \& Raquel Picazo}
Queridos Delegados,

¡Bienvenidos a la séptima edición del modelo de Naciones Unidas del IE University, IEUMUN 2026! Nosotras somos Jimena Solórzano y Raquel Picazo, y juntas formamos la presidencia de este comité. Es un honor participar en esta edición de la conferencia como moderadoras del Fondo de las Naciones Unidas para la Infancia y tener la oportunidad de participar en estos días de debate de los cuales vosotros, delegados, seréis protagonistas. Ya sea vuestro primer, segundo u octavo modelo, esperamos que durante el transcurso de estos días os divirtáis, aprendáis algo nuevo, disfrutéis de conocer nuevas personas y, sobre todo, que nos involucremos en debates divertidos a la vez que fructuosos. Antes de comenzar, nos gustaría presentarnos y que nos conozcáis un poco mejor.

Hola de nuevo, mi nombre es Jimena y seré vuestra presidenta de comité. Estoy en mi tercer año de Economía y Relaciones Internacionales en el CEU San Pablo.

Originalmente soy de Costa Rica, pero me crié en República Dominicana. Gracias a mi carrera y a mi experiencia de vida he tenido la oportunidad de exponerse a diferentes perspectivas políticas, y espero que a lo largo de este debate pueda escuchar ideas innovadoras de vuestra parte. Al tener experiencia tanto como delegada como en presidencia, estoy muy emocionada de ver qué propuestas vais a traer al comité; no dudo en lo absoluto en que estaréis al nivel de este modelo. Y desde luego, si tenéis alguna duda, tanto previa como durante el debate, no dudéis en acercaros a nosotras, que estamos más que dispuestas a ayudaros. ¡No veo la hora de conoceros, nos vemos en noviembre!

¡Hola! Me llamo Raquel y seré parte de la presidencia de este comité junto a Jimena. Soy de Madrid y actualmente estoy en segundo año en IE University estudiando un doble grado en Derecho y Relaciones Internacionales. A pesar de estar en el inicio de mi carrera, mantengo un interés creciente en las instituciones internacionales y siempre me ha llamado la atención la ONU especialmente tratando desarrollo, paz y seguridad. A lo largo de mis cinco años involucrada en Modelos de Naciones Unidas he tenido la oportunidad de conocer gente maravillosa que a día de hoy forman parte de mi vida así como de aprender mediante el debate sobre los retos a los que se enfrenta el mundo día a día. Espero que este comité os presente un reto divertido y tengáis la oportunidad de aprender y de escuchar opiniones desde el mayor número de puntos de vista posibles a lo largo de la conferencia. ¡Estoy súper ilusionada de ver qué retos, soluciones e innovaciones nos presentáis!

De nuevo, esperamos que disfrutéis de la conferencia lo máximo posible. Como moderadoras, estamos muy entusiasmadas de ver qué nos presentáis durante estos días en cuanto al tema “Protección de los menores frente a la radicalización y explotación digital”. Esperamos que, mediante nuestra guía y vuestras maravillosas ideas, nos encontremos en un ecosistema de creatividad, innovación y, por supuesto, confianza que os impulse a ofrecer lo mejor de vosotros mismos, así como haremos nosotras.

¡Bienvenidos al Fondo de las Naciones Unidas para la Infancia! Jimena Solórzano y Raquel Picazo.

\end{chairletter}
